\section{Applications of curves cut out by quadrics}
\label{app:applications}

We prove the cases of Theorem~\ref{intro:main} that can be obtained from
complete intersection curves of quadrics. All capacities in this section are
normalized to a target of capacity one. We use the surface packing
Lemma~\ref{surf:packing} and the geometric constructions of the preceding
section with their neighborhood and relative avoidance conclusions.

\subsection{Complete intersections and their normal models}

\begin{lemma}\label{app:quadrics}
Let $N\ge3$, and let $F\subset\PP^N$ be a hyperplane. There is a smooth
connected curve $C$, the complete intersection of $N-1$ quadrics, transverse
to $F$. If $A=\cO_{\PP^N}(1)|_C$, then
\begin{equation}\label{app:quadric-data}
 \deg A=2^{N-1},\qquad
 N_{C/\PP^N}=(A^{\otimes2})^{\oplus(N-1)},\qquad
 g(C)=1+(N-3)2^{N-2}.
\end{equation}
Writing $H$ for the hyperplane pullback and $E_C$ for the exceptional divisor
of $\Bl_C\PP^N$, the rational class $H-dE_C$ is ample whenever
$0<d<1/2$. The class
\[
 A+dU\quad\hbox{on}\quad
 \PP_C\bigl(\mathbf1\oplus (A^{\otimes2})^{\oplus(N-1)}\bigr)
\]
is ample for the same range of rational $d$.
\end{lemma}

\begin{proof}
The quadratic Veronese system restricts to a generated linear system
defining a closed immersion on every smooth intermediate intersection,
and on its intersection with $F$. Bertini's theorem
\cite[Tag~0FD6]{Stacks}, applied successively to both, therefore gives
$N-1$ quadrics whose intersection is a smooth curve transverse to $F$.
The Koszul
resolution of its structure sheaf has terms that are sums of
$\cO_{\PP^N}(-2j)$, $1\le j\le N-1$. The vanishings
$H^i(\PP^N,\cO(t))=0$ for $0<i<N$
\cite[Tag~01XT]{Stacks}, together with the absence of global
sections of these negative twists, give
$H^0(C,\cO_C)=\C$. In particular, $C$ is connected.

The intersection class is $(2H)^{N-1}$, so intersection with one further
hyperplane gives $\deg A=2^{N-1}$. The defining equations form a regular
sequence; hence their classes identify the conormal bundle with
$(A^{-2})^{\oplus(N-1)}$. This proves the assertion about the normal bundle.
Taking determinants in the tangent-normal exact sequence gives
\[
 K_C=K_{\PP^N}|_C\otimes\det N_{C/\PP^N}
     =A^{\otimes(N-3)}.
\]
Consequently $2g(C)-2=(N-3)2^{N-1}$, proving
\eqref{app:quadric-data}. These uses of Bertini, the Koszul resolution,
and adjunction are in their usual smooth projective form
\cite{Hartshorne1977}; the regular-sequence and conormal statements
are also given in \cite[Tags~063C and~06AA]{Stacks}.

The quadric equations generate $\mathcal I_C(2)$, and hence give a
surjection of graded sheaves of algebras
\[
 \cO_{\PP^N}[T_1,\ldots,T_{N-1}]
 \longrightarrow\bigoplus_{j\ge0}\mathcal I_C^j(2j).
\]
The relative Proj of the target is $\Bl_C\PP^N$: twisting the degree-$j$
part of the Rees algebra by $\cO(2j)$ leaves its relative Proj unchanged.
Thus this surjection gives a closed embedding
\[
 \Bl_C\PP^N\hookrightarrow\PP^N\times\PP^{N-2}.
\]
The hyperplane classes of the factors pull back to $H$ and $2H-E_C$.
For rational $0<d<1/2$, the identity
\[
 H-dE_C=(1-2d)H+d(2H-E_C)
\]
expresses this class as the restriction of a positive rational product
polarization. After clearing denominators, a Segre--Veronese embedding
shows that it is ample.

Finally, choose a positive curvature form $\sigma$ for $A$.
Lemma~\ref{geo:bundle}, with normal exponent two, gives a K\"ahler form in
$A+dU$: its horizontal coefficient is $1-2\sum p_j\ge1-2d>0$, also at
infinity. A rational class admitting a K\"ahler representative is ample
by the Kodaira criterion, after clearing denominators.
\end{proof}

Here and below an angular connection on a line bundle has period one.
For an integer $m\ge1$, write $t=\sum_{j=1}^m p_j$. Integration over the
standard simplex gives
\begin{equation}\label{app:simplex-integrals}
 \int_{\{p\ge0:t\le h\}}f(t)\,dp
 =\frac1{(m-1)!}\int_0^h f(t)t^{m-1}\,dt,
 \qquad
 \int_{\R_{\ge0}^m}(r-t)_+\,dp=\frac{r^{m+1}}{(m+1)!}.
\end{equation}
For the first identity, the volume of $\{p\ge0:t\le h\}$ is $h^m/m!$,
by scaling the unit simplex, whose volume is $1/m!$ by iterated integration.
Differentiating this volume gives the first formula, initially for
continuous $f$. The second follows by taking $f(t)=r-t$ on $[0,r]$.

\subsection{All capacities below one half}

\begin{proposition}\label{app:small}
Let $n\ge3$, and let $\rho_1,\ldots,\rho_k$ be positive real numbers with
\[
 \rho_i<\tfrac12\quad(1\le i\le k),\qquad
 \sum_{i=1}^k\rho_i^n<1.
\]
Then the closed balls of capacities $\rho_i$ admit a disjoint symplectic
packing into $\Int B^{2n}(1)$, with embeddings defined on neighborhoods
of the closed source balls.
\end{proposition}

\begin{proof}
Set $m=n-1$. The displayed inequalities are strict and there are finitely
many of them. By continuity and density of the rationals, choose rational
numbers $r_i>\rho_i$ such that
\begin{equation}\label{app:small-slack}
 r_*:=\max_i r_i<\tfrac12,\qquad \sum_i r_i^n<1.
\end{equation}
For $0<h<1/2$, formula~\eqref{app:simplex-integrals} gives
\begin{equation}\label{app:small-volume}
 J_m(h):=\int_{\{t\le h\}}2^m(1-2t)\,dp
 =\frac{2^mh^m}{m!}\left(1-\frac{2mh}{m+1}\right)
 \longrightarrow\frac1{n!}
 \quad(h\uparrow\tfrac12).
\end{equation}
It follows from \eqref{app:small-slack} that we can choose rational
parameters satisfying
\begin{equation}\label{app:small-truncation}
 r_*<h<d<\tfrac12,\qquad
 J_m(h)>\frac1{n!}\sum_i r_i^n.
\end{equation}
Indeed, choose $h$ sufficiently close to $1/2$ and then choose $d$ strictly
between $h$ and $1/2$; each requirement is open.

Apply Lemma~\ref{app:quadrics} in $\PP^n$, with the forbidden hyperplane
$H_\infty$. Deform to the normal cone of the resulting curve $C$, using
the class $H-d[Y]$. On the two central components the polarizations are
$H-dE_C$ and $A+dU$, respectively. Both are ample by that lemma, so
Proposition~\ref{geo:transfer} applies. On the normal component use the
form of Lemma~\ref{geo:bundle}. Its affine part, truncated at $t\le h$, is
\begin{equation}\label{app:small-model}
 \sigma+\dd\sum_{j=1}^m p_j\alpha_j,
 \qquad \dd\alpha_j=-2\sigma,\qquad
 \int_C\sigma=2^m.
\end{equation}
The strict inequality $h<d$ supplies a neighborhood of the entire
truncated region in the affine part.

Let $C^o=C\setminus(C\cap H_\infty)$. The removed set is finite, and $C^o$
has the positive genus calculated in \eqref{app:quadric-data}. Exhaust it
by compact connected smooth bordered surfaces $\Sigma_\ell$, with collars
contained in $C^o$. For example, remove progressively smaller disjoint
coordinate disks about the punctures. Each such surface has the same genus
as $C$. Since $H^2(\Sigma_\ell;\Z)=0$, the normal line bundles admit smooth
unitary frames on a slightly larger bordered surface. Write in those frames
$\alpha_j=\dd\theta_j+A_j$, where $\dd A_j=-2\sigma$, and choose
$\gamma_\ell$ with $\dd\gamma_\ell=\sigma$. The form
\eqref{app:small-model} becomes
\[
 \omega_V+\dd\Gamma_{0,\ell},\qquad
 \Gamma_{0,\ell}=\gamma_\ell+\sum_jp_jA_j,\qquad
 \dd_C\Gamma_{0,\ell}=(1-2t)\sigma>0
\]
on a neighborhood of $\Sigma_\ell\times V$, for
$\Delta=\{p\ge0:t\le h\}$. The polar notation is smooth at every axis:
if $z_j=x_j+iy_j$, then
$p_j\dd\theta_j=(x_j\dd y_j-y_j\dd x_j)/2$.
The integrated horizontal forms converge uniformly on $\Delta$ to
\[
 H_0(p)=2^m(1-2t),
\]
since $\int_{\Sigma_\ell}\sigma\to\int_C\sigma$ and $1-2t$ is bounded
there.

We verify all the remaining hypotheses of the surface packing lemma.
Each auxiliary simplex $\{t\le r_i\}$ lies strictly inside the outer face
of $\Delta$, by \eqref{app:small-truncation}. The function
\[
 P(p)=H_0(p)-\sum_i(r_i-t)_+
\]
is continuous and concave, being affine minus a sum of convex functions.
Choose $r_*/h<\tau<1$. On $\Delta\setminus\tau\Delta$ every cap vanishes,
and hence
\[
 P(p)=H_0(p)\ge2^m(1-2h)>0.
\]
Finally, \eqref{app:simplex-integrals} and
\eqref{app:small-truncation} give
\[
 \int_\Delta P\,dp=J_m(h)-\frac1{n!}\sum_i r_i^n>0.
\]
Thus $P$ has positive mean and meets the hypotheses of
Lemma~\ref{cmp:comparison}. All the hypotheses of
Lemma~\ref{surf:packing} now hold. Since $\rho_i<r_i$, that lemma provides
a compact packing of the requested closed balls in the interior of a
sufficiently large surface model.

This packing lies away from the infinity locus of the normal component
and from its fibers over $C\cap H_\infty$. These forbidden fibers are
smooth and pairwise disjoint. Proposition~\ref{geo:transfer} therefore
transfers the packing into $\PP^n\setminus H_\infty$ with the standard
unit Fubini--Study form; the prescribed-form conclusion uses
Lemma~\ref{geo:relative-moser} preserving $H_\infty$. The usual projective
moment coordinates identify this complement with $\Int B^{2n}(1)$.
All transfers are defined on neighborhoods of the compact packing.
Together with $h<d$ and the capacity slack $\rho_i<r_i$, this proves the
asserted neighborhood statement.
\end{proof}

\subsection{A distinguished large ball in dimension at least eight}

\begin{proposition}\label{app:large}
Let $n\ge4$, let $A\ge1/2$, and let $\rho_1,\ldots,\rho_k>0$ satisfy
\[
 A^n+\sum_i\rho_i^n<1,\qquad A+\rho_i<1\quad(1\le i\le k).
\]
Then the closed ball of capacity $A$ and the closed balls of capacities
$\rho_i$ admit a disjoint symplectic packing into $\Int B^{2n}(1)$, with
embeddings on neighborhoods of all source balls.
\end{proposition}

\begin{proof}
The case of no remaining balls is immediate, so suppose $k\ge1$.
Continuity of the finitely many strict inequalities allows rational
parameters $a,r_1,\ldots,r_k$ with
\begin{equation}\label{app:large-slack}
 \tfrac12<a<1,\qquad A<a,\qquad \rho_i<r_i<s:=1-a,
 \qquad\sum_i r_i^n<1-a^n.
\end{equation}
For completeness, the point $(A,\rho_1,\ldots,\rho_k)$ has a neighborhood
on which the volume and pairwise inequalities remain strict. Choose $a>A$
in that neighborhood, also $a>1/2$, and then choose the $r_i>\rho_i$ within
it. Rational choices exist because all increments may be taken arbitrarily
small. In particular $0<s<1/2$.

Lemma~\ref{geo:shell} lets us reserve the central ball of capacity $a$
and seek the remaining packing in its exterior shell, whose volume is
$(1-a^n)/n!$. We construct surface models of this shell in two stages:
first a normal line over a hyperplane, then a normal model of a curve
inside that hyperplane. The second model will be embedded symplectically
in the base of the first; pulling back the first normal line then gives
a model over the curve.

Put $q=n-2$ and $r_*:=\max_i r_i$. For now take rational parameters with
\[
 r_*<h_1<d_1<s,\qquad r_*<h_2<d_2<\tfrac12.
\]
The construction below works for every such choice and every sufficiently
large bordered piece of the curve. We will fix the parameters after
computing the model's volume.

\paragraph{The two normal constructions.}
Take $T=\Bl_o\PP^n$, $L=H-aE$, and the two disjoint forbidden divisors
$E,H_\infty$ of Lemma~\ref{geo:shell}. Choose a hyperplane not containing
$o$ and different from $H_\infty$; its strict transform $D$ is isomorphic
to $\PP^{n-1}$ and is disjoint from $E$. Its normal line is
$\cO_D(1)$, and $L|_D=\cO_D(1)$. Choose it so that
$F=D\cap H_\infty$ is a hyperplane in $D$.

Deform to the normal cone of $D$ with parameter $d_1$. The strict component
has class $(1-d_1)H-aE$, which is ample: on the blowup of projective space at
a point, $xH-yE$ is ample for $x>y>0$, and here $1-d_1>a>0$.
This ampleness criterion is proved in Lemma~\ref{geo:shell}.
The other component is
\[
 Y_1=\PP_D(\mathbf1\oplus\cO_D(1)),\qquad
 \cO_D(1)+d_1U_1.
\]
Fix a unit Fubini--Study form $\omega_D$ on $D$, and write $u$ for
the moment of the first normal line. By Lemma~\ref{geo:bundle}, this
class is ample and has an affine normal model
\begin{equation}\label{app:first-model}
 \omega_D+\dd(u\alpha_1),\qquad
 \dd\alpha_1=-\omega_D,\qquad 0\le u<d_1.
\end{equation}
There is a neighborhood of the truncation $u\le h_1$ in this model.
The forbidden preimage in $Y_1$ lies over $F$; the preimage of $E$ is empty.

Apply Lemma~\ref{app:quadrics} inside the fixed base $D=\PP^{n-1}$, now with
$N=n-1$ and forbidden hyperplane $F$. We obtain a curve $C$ of degree $2^q$
and genus
\[
 g(C)=1+(n-4)2^{n-3}\ge1,
\]
with normal bundle $\cO_C(2)^{\oplus q}$. The second normal construction,
with parameter $d_2$, is ample on both components by
Lemma~\ref{app:quadrics}. Write $\sigma$ for a positive curvature form of
$\cO_C(1)$; its integral is $2^q$. Its affine normal model is
\begin{equation}\label{app:second-model}
 \omega_X=\sigma+\dd\lambda,\qquad
 \lambda=\sum_{j=2}^{n-1}v_j\alpha_j,\qquad
 \dd\alpha_j=-2\sigma,
 \qquad \sum_{j=2}^{n-1}v_j<d_2.
\end{equation}

Exhaust $C\setminus F$ by bordered surfaces $\Sigma$, as in the proof of
Proposition~\ref{app:small}. For each such $\Sigma$, the truncation
$\sum v_j\le h_2$ is a compact region away from infinity and from the
finitely many forbidden fibers. Proposition~\ref{geo:transfer}, with its
prescribed K\"ahler form conclusion, gives a symplectic embedding of a
neighborhood of this region into $(D\setminus F,\omega_D)$. After shrinking that neighborhood, choose an open collar enlargement
$\Sigma^+$ of $\Sigma$, with compact closure in $C\setminus F$ and
deformation retracting onto $\Sigma$. We can take the domain $X$ to be an
open disk bundle over $\Sigma^+$ with fiberwise star-shaped fibers,
containing a neighborhood of the prescribed compact truncation.
Denote this embedding by
\[
 \phi:(X,\omega_X)\longrightarrow(D\setminus F,\omega_D).
\]
All line bundles over the bordered curve pieces will be smoothly framed;
the embeddings $\phi$ are only required to be symplectic.

\paragraph{Lifting the base embedding and changing coordinates.}
Pull back the Hermitian normal line in \eqref{app:first-model}, with its
unitary connection, along $\phi$. Its angular connection has curvature
$-\phi^*\omega_D=-\omega_X$. Let $\operatorname{pr}:X\to\Sigma^+$ denote the bundle
projection and $i:\Sigma^+\to X$ its zero section. The contraction of the
disk fibers gives $H^2(X;\Z)=H^2(\Sigma^+;\Z)=H^2(\Sigma;\Z)=0$. Hence the pulled-back line
and $\operatorname{pr}^*\cO_C(1)$ admit a smooth unitary bundle identification.
Let $\alpha_{1,C}$ be the angular connection on the latter line, pulled
back from the curve and with curvature $-\sigma$. The connection
\[
 \alpha_{1,C}-\lambda
\]
has curvature $-\sigma-\dd\lambda=-\omega_X$ too. The difference between
the pulled-back connection and this reference connection is therefore a
closed ordinary one-form $\xi$ on $X$.

Set $\eta=i^*\xi$, which is a closed one-form on $\Sigma^+$. Radial contraction
in the disk fibers and the homotopy formula give
\[
 \xi=\operatorname{pr}^*\eta+\dd f
\]
for a smooth real function $f$ on $X$. For instance, in a unitary
trivialization the contraction is $(x,z)\mapsto(x,tz)$, and one can take
$f(x,z)=\int_0^1\xi_{(x,tz)}(0,z)\,dt$. A unitary gauge change removes
$\dd f$. Thus, omitting pullback symbols, the pulled-back first connection
has the precise expression
\begin{equation}\label{app:lifted-connection}
 \alpha_{1,C}-\sum_{j=2}^{n-1}v_j\alpha_j+\eta.
\end{equation}
There is no assertion that $\eta$ is exact; its periods are retained.
These operations are smooth unitary operations, so they preserve the first
fiber moment $u$ and require no holomorphic lift of $\phi$.

The lifted form \eqref{app:first-model} is consequently
\begin{equation}\label{app:lifted-form}
 \sigma+\dd\left(u(\alpha_{1,C}+\eta)
                  +(1-u)\sum_{j=2}^{n-1}v_j\alpha_j\right).
\end{equation}
The terms $v_j\alpha_j$ extend smoothly over zero: in a unitary frame,
$\alpha_j=\dd\theta_j+A_j$ and
$v_j\dd\theta_j=(x_j\dd y_j-y_j\dd x_j)/2$.
Use the disk coordinates $z_1,z_2,\ldots,z_{n-1}$ with
$u=\pi|z_1|^2$ and $v_j=\pi|z_j|^2$. The actual smooth change of coordinates
\begin{equation}\label{app:moment-change}
 \zeta_1=z_1,\qquad
 \zeta_j=\sqrt{1-\pi|z_1|^2}\,z_j\quad(2\le j\le n-1)
\end{equation}
is a diffeomorphism on $u<1$, with smooth inverse there, including all
coordinate axes. Its moments are
\[
 p_1=u,\qquad p_j=(1-u)v_j\quad(j\ge2),
\]
and its angles agree with the previous ones. Write
$w=\sum_{j=2}^{n-1}p_j$. The truncations $u\le h_1$ and
$\sum_{j=2}^{n-1}v_j\le h_2$ become
\begin{equation}\label{app:large-polytope}
 \Delta=\Delta(h_1,h_2)=
 \{p\ge0:u\le h_1,\ w+h_2u\le h_2\}.
\end{equation}

Writing $\alpha_{1,C}=\dd\theta_1+A_1$ with $\dd A_1=-\sigma$, and choosing
$\gamma$ with $\dd\gamma=\sigma$, expression \eqref{app:lifted-form} becomes
\begin{equation}\label{app:large-surface-form}
 \omega_V+\dd\Gamma_0,\qquad
 \Gamma_0=\gamma+p_1(A_1+\eta)+\sum_{j=2}^{n-1}p_jA_j,\qquad
 \dd_C\Gamma_0=(1-u-2w)\sigma.
\end{equation}
This is a horizontal smooth toric primitive in precisely the coordinates
used by Lemma~\ref{surf:packing}. Its horizontal coefficient is bounded
below on $\Delta$ by
\begin{equation}\label{app:large-positive}
 1-u-2w\ge(1-u)(1-2h_2)
             \ge(1-h_1)(1-2h_2)>0.
\end{equation}
The strict truncation inequalities, the base collars, and the smooth
inverse in \eqref{app:moment-change} give the model on a neighborhood of
$\Sigma\times V$. For each fixed choice of the truncation parameters, the constructions
can be made for arbitrarily large $\Sigma$. Their integrated horizontal
coefficients converge uniformly to
\[
 H_0(p)=2^q(1-u-2w),
\]
because $\int_\Sigma\sigma\to2^q$ and $1-u-2w$ is bounded on the fixed
$\Delta$.

\paragraph{Choosing a model with enough volume.}
For $0<h_2<1/2$, integration in the $q$ moments different from $u$ gives
\[
 \int_{\{w\le h_2(1-u)\}}2^q(1-u-2w)\,dp_2\cdots dp_{n-1}
 =B_q(h_2)(1-u)^{q+1},
\]
where
\[
 B_q(h_2)=2^q\left(\frac{h_2^q}{q!}
             -\frac{2q h_2^{q+1}}{(q+1)!}\right)
 \longrightarrow\frac1{(q+1)!}\quad(h_2\uparrow\tfrac12).
\]
Thus
\begin{equation}\label{app:large-volume}
 \int_{\Delta(h_1,h_2)}H_0\,dp
 =B_q(h_2)\frac{1-(1-h_1)^n}{n}
 \longrightarrow\frac{1-a^n}{n!}
 \quad\bigl(h_1\uparrow s,\ h_2\uparrow\tfrac12\bigr).
\end{equation}
Now fix these parameters. By \eqref{app:large-slack}, they can be
chosen rationally so that
\begin{equation}\label{app:large-truncation}
 \begin{gathered}
 r_*<h_1<d_1<s,\qquad
 r_*<h_2<d_2<\tfrac12,\\
 \int_{\Delta(h_1,h_2)}H_0\,dp>
                 \frac1{n!}\sum_i r_i^n.
 \end{gathered}
\end{equation}
Indeed, the last inequality holds for $h_1,h_2$ sufficiently close to their
limits, where the first two lower bounds also hold; $d_1,d_2$ can then be
inserted between the chosen truncations and their respective limits.

\paragraph{The surface packing conditions and transfer to the shell.}
The defining outer inequalities of $\Delta$ have nonnegative coefficients,
as required by Lemma~\ref{cmp:comparison}. If $u+w\le r_i$, then
\[
 u\le r_i<h_1,\qquad
 w+h_2u\le w+u\le r_i<h_2.
\]
Thus every auxiliary simplex is contained in $\Delta$ strictly away from
its outer faces. The difference
\[
 P(p)=H_0(p)-\sum_i(r_i-u-w)_+
\]
is continuous and concave. Choose
\[
 \max\{r_*/h_1,r_*/h_2\}<\tau<1.
\]
The same inequalities show that all cap supports lie in $\tau\Delta$,
whose outer equations are $u\le\tau h_1$ and
$w+h_2u\le\tau h_2$. Outside that smaller polytope, $P=H_0$, and
\eqref{app:large-positive} gives the uniform positive lower bound
$2^q(1-h_1)(1-2h_2)$. Finally,
\eqref{app:simplex-integrals} and \eqref{app:large-truncation} give
\[
 \int_\Delta P\,dp
 =\int_\Delta H_0\,dp-\frac1{n!}\sum_i r_i^n>0.
\]
These verifications establish the hypotheses of
Lemmas~\ref{cmp:comparison} and~\ref{surf:packing}.

The latter lemma packs the closed balls of capacities $\rho_i<r_i$ in one
of the sufficiently large models \eqref{app:large-surface-form}. Composing
with the smooth coordinate identification and the lifted embedding puts
this compact packing in $Y_1$, with $u<d_1$, over $D\setminus F$.
It avoids both infinity and the forbidden preimage, which is the smooth
projective line bundle over $F$. In the first normal construction,
$E$ and $H_\infty$ are disjoint smooth divisors; their intersections with
$D$ are respectively empty and $F$. Proposition~\ref{geo:transfer}
therefore transfers the compact packing to
$T\setminus(E\cup H_\infty)$ with a K\"ahler form in $H-aE$.
Lemma~\ref{geo:shell} identifies this complement symplectically with
\[
 \{z\in\C^n:a<\pi\textstyle\sum_j|z_j|^2<1\}.
\]
The distinguished closed ball of capacity $A<a$ occupies the central
standard ball and is disjoint from this shell packing. All constructed
images are compact subsets of the open target. Every embedding used above
is defined on a neighborhood of the relevant compact subset; shrinking
those neighborhoods under the finitely many compositions retains the
neighborhood embeddings of the source balls and their disjointness.
\end{proof}
